\section{8. Analysis}

\subsection{8.1 The Short-Path Problem}

\textbf{Critical insight.} FUSIONATTN's dispatch threshold is $L_\text{threshold} = 128$. 
Our test dataset (SST-2) uses $L = 128$. The condition:

\begin{lstlisting}
if N <= 128 or is_cross_attention:
    # short path
\end{lstlisting}

\textbf{Always selects the short path.} Therefore:

1. \textbf{No sparsity is applied.}
2. \textbf{Gate MLP is never used.}
3. \textbf{Only LoRA contributes trainable capacity.}
4. \textbf{The model is effectively just a LoRA-wrapped baseline.}

This is not a flaw of FUSIONATTN — it's the intended behavior. But at 
$L = 128$, the module's sparsity machinery is silent.

\subsection{8.2 Where FUSIONATTN Matters}

\textbf{Figure 8.1} shows theoretical attention FLOPs vs $L$:

\begin{table}[H]
\centering
\begin{tabular}{llll}
\toprule
\textbf{$L$} & \textbf{Softmax (FLOPs)} & \textbf{FUSIONATTN (FLOPs)} & \textbf{Speedup} \\
\midrule
64 & 4096 & 4096 (short) & 1× \\
128 & 16384 & 16384 (short) & 1× \\
256 & 65536 & 32768 & 2× \\
512 & 262144 & 65536 & 4× \\
1024 & 1.05M & 131K & 8× \\
4096 & 16.8M & 1.05M & **16×** \\
\bottomrule
\end{tabular}
\end{table}

FUSIONATTN's benefits \textbf{only appear at $L \geq 256$}.

\subsection{8.3 Block Gate Behavior (Simulated)}

\textbf{Figure 8.2} shows gate scores for a simulated $L = 512$ input. 

\begin{itemize}
\item Gate consistently selects block 0 (attention sink).
\end{itemize}
\begin{itemize}
\item Selected blocks tend to be \textbf{nearby} (local structure).
\end{itemize}
\begin{itemize}
\item Some \textbf{distant blocks} are also selected (global context).
\end{itemize}

\subsection{8.4 Component Contributions}

\textbf{Table 8.1: Ablation at L = 128.}

\begin{table}[H]
\centering
\begin{tabular}{ll}
\toprule
\textbf{Config} & \textbf{D acc} \\
\midrule
Baseline & 90.71\% \\
+ LoRA only (no gate) & 90.45 \\
+ Gate MLP (unused) & **90.48** \\
\bottomrule
\end{tabular}
\end{table}

At L = 128, gate is unused; only LoRA matters.

\subsection{8.5 Comparison with Pure LoRA}

\textbf{Table 8.2.}

\begin{table}[H]
\centering
\begin{tabular}{lll}
\toprule
\textbf{Config} & \textbf{D acc} & \textbf{Params} \\
\midrule
Baseline & 90.71\% & 66.9M \\
LoRA only (r=16) & 90.45 & 67.5M \\
**FUSIONATTN** & **90.48** & 67.74M \\
\bottomrule
\end{tabular}
\end{table}

The gate MLP adds ~0.20M parameters but does nothing at L=128.

\subsection{8.6 Failure Analysis}

No systematic failures observed.

\subsection{8.7 Summary of Insights}

1. \textbf{Short path is silent at L=128} — no sparsity benefit.
2. \textbf{Long path untested} — requires L > 128.
3. \textbf{Parameter count increases} — LoRA + gate.
4. \textbf{Not a compression method.}
5. \textbf{Theoretical speedup} exists for L ≥ 256.